\documentclass[11pt,a4paper]{article}
\usepackage[T1]{fontenc}
\usepackage{lmodern}
\usepackage[margin=27mm]{geometry}
\usepackage{amsmath,amssymb,amsthm,booktabs,microtype}
\usepackage[hidelinks]{hyperref}
\hypersetup{pdftitle={Useful Constraints or Compression by Group Counts},
  pdfsubject={Max 2-SAT branching and structural algorithms}}
\newtheorem{theorem}{Theorem}
\newtheorem{proposition}[theorem]{Proposition}
\newtheorem{corollary}[theorem]{Corollary}
\theoremstyle{definition}
\newtheorem*{desiredtheorem}{Desired theorem (unproved)}
\newcommand{\Ostar}{O^{*}}
\newcommand{\ind}[1]{\mathbf 1\{#1\}}
\emergencystretch=2em
\clubpenalty=10000
\widowpenalty=10000
\title{Useful Constraints or Compression by Group Counts\\
\large Literal neighborhoods, repeated interactions, and adaptive branching}
\author{}
\date{10 October 2026}
\begin{document}
\maketitle

\begin{abstract}
The branching algorithm chooses which constraints to create and when to
create them. A family of weak constraints that it can decline to add is
therefore not an obstruction to the program. Likewise, redundant information
that admits a faster exact solver belongs on the structural side of the
desired dichotomy. We make this principle concrete. First, we derive the
one-literal-neighborhood pruning rule using the occurrence gap, including
unit clauses. Second, we prove a polynomial-space algorithm for a cut whose
rows fall into repeated signed interaction profiles, allowing arbitrary
interactions within either side. Its saving is measured by
$\Gamma=\sum_i(|A_i|-\log_2(|A_i|+1))$. Third, branching on exceptional
columns can merge groups; an explicit cost bound accounts for every such
branch and gives a variable-selection score. Literal bicliques and the
earlier aligned example are special cases. The missing result is an
availability theorem connecting failure of useful constraint generation to
an exploitable structure, with a bound on the total search cost.
\end{abstract}

\section{The algorithmic question}

We seek an exact deterministic algorithm with polynomial space and time
$\Ostar(c^n)$, for a fixed $c<2$ independent of a density bound
$\rho(n)=m_2/n$ tending to infinity. The eventual target is the full
$\rho=O(n)$ range for canonical simple formulas. Unless stated otherwise,
there is at most one binary clause per unordered pair of distinct variables;
unit clauses may have multiplicity. The group-count theorem also holds for
general rational quadratic coefficients of polynomial encoding length.

The intended dichotomy is algorithmic: create constraints that give useful
\emph{additional} pruning, expose variables so that useful constraints become
available, or exploit the structure responsible for redundant information.
There is no requirement to add every available constraint. Useful pruning
must be assessed relative to the current domain or an inherited completion
certificate, rather than by counting supports alone. A structural solver is
allowed to avoid enumerating a large feasible domain altogether.

The results below are self-contained derivations in this research program.
The local algorithms and identities are proved; the global availability
statement in Section~\ref{sec:target} is a desired theorem. No claim of
literature novelty is made. The relation to the existing notes is recorded
in Section~\ref{sec:relation}.

\section{Pruning from just one literal neighborhood}\label{sec:one}

For a literal $\ell$, let $N(\ell)$ contain the partner literal $m$ for
each binary clause $(\ell\lor m)$. Define the occurrence gap
\begin{equation}
 G_\ell=d_2(\ell)+d_1(\ell)-d_2(\neg\ell)-d_1(\neg\ell).
 \label{eq:gap}
\end{equation}
Unit and binary occurrences count equally in this definition. For a fixed
assignment $y$ to all other variables, put
\[
 A_\ell(y)=\sum_{m\in N(\ell)}\ind{m\text{ is true}},\qquad
 B_\ell(y)=\sum_{m\in N(\neg\ell)}\ind{m\text{ is true}}.
\]

\begin{proposition}[One-neighborhood necessary condition]\label{prop:one}
The exact flip difference is
\begin{equation}
 F(\ell=0,y)-F(\ell=1,y)=A_\ell(y)-B_\ell(y)-G_\ell.
 \label{eq:flip}
\end{equation}
Consequently, the strict-improvement condition in the second branch
implies the sound one-neighborhood pruning constraint
\begin{equation}
 \boxed{A_\ell(y)\ge G_\ell+1.} \label{eq:one}
\end{equation}
Here the first branch is $\ell=1$, and the constraint inherits the scope
and tie convention of the established flip-based pruning rule.
\end{proposition}
\begin{proof}
Switching from $\ell=1$ to $\ell=0$ changes the score of the positive
binary clauses by $A_\ell-d_2(\ell)$, of the negative binary clauses by
$d_2(\neg\ell)-B_\ell$, and of the units by
$d_1(\neg\ell)-d_1(\ell)$. Their sum gives~\eqref{eq:flip}.
An integer-valued strict improvement requires
$A_\ell-B_\ell\ge G_\ell+1$; use $B_\ell\ge0$.
\end{proof}

Thus the exact comparison combines both polarities, but a valid derived
constraint can involve only the neighbors of one literal. Its support is
directly a literal neighborhood. For example, when $G_\ell=0$,
\eqref{eq:one} is the ordinary clause $\bigvee_{m\in N(\ell)}m$.
The algorithm can retain the combined inequality, a one-sided consequence,
or both, depending on the information and representation costs. A
one-sided consequence is necessary for strict improvement, not sufficient;
its implication by older constraints does not by itself certify that the
second value of the branching variable is better.

The combined form agrees with the earlier deficit notation. With
$r=d_2(\ell)$ and
$p_\ell=d_2(\neg\ell)-d_1(\ell)+d_1(\neg\ell)$,
\[
 (r-A_\ell)+B_\ell\le p_\ell-1.
\]
Dropping either nonnegative term yields a valid one-neighborhood bound.
One must not split $p_\ell-1$ into independently imposed smaller budgets.
These statements use the existing optimization-preserving pruning
convention, rather than assuming that arbitrary hard constraints are closed
under a variable flip.

\section{Literal bicliques and the meaning of overlap}\label{sec:biclique}

A literal biclique on disjoint variable sets $A,B$ specifies one literal
per variable and contains the clause $(\ell_u\lor m_v)$ for every
$u\in A,v\in B$. Write $s=|A|$, $t=|B|$. Given such a biclique, assign
$R=V\setminus(A\cup B)$. Complement variables so that all specified
literals are positive. Among assignments to $A$, retain the best score of
the non-cross terms for each number $i$ of true variables; do the same for
$B$ and $j$. The cross score is
\[
 st-(s-i)(t-j).
\]
Enumerating the two sides separately and combining the count tables gives
\begin{equation}
 \Ostar\bigl(2^{n-s-t}(2^s+2^t)\bigr)
 =\Ostar\bigl(2^{n-\min\{s,t\}}\bigr)
 \label{eq:biclique}
\end{equation}
time and polynomial space. All other clauses and all units are arbitrary.
The sets are supplied or identified; this calculation does not pay for an
unrestricted search over all partitions.

For fixed positive $\alpha,\beta$, absence of $K_{\alpha n,\beta n}$
implies that every $\alpha n$ literal neighborhoods have fewer than
$\beta n$ common literals. This is a statement about intersections.
It is relevant to the supports in~\eqref{eq:one}, but a proof of the
branching strategy must also use the thresholds of the constraints
\emph{actually selected}. The absence condition alone is not an
availability or running-time theorem.

There is a simple robust version of~\eqref{eq:biclique}. If every literal
on an $s$-element side misses at most a fixed integer $b\ge1$ members of
a $t$-element opposite side, take
$s'=\min\{s,\lfloor t/(2b)\rfloor\}$ centers. Delete the union of their
at most $bs'$ missing partners. At least $t/2$ opposite literals remain,
giving a complete biclique with both sides linear whenever $s,t$ are
linear. This is one way to resolve nearly repeated neighborhoods directly.

\section{A redundant family that should be simplified}\label{sec:aligned}

For each variable pair independently choose either $(u\lor v)$ or
$(\neg u\lor\neg v)$, with equal probability. For fixed
$\alpha,\beta>0$ with $\alpha+\beta\le1$, this family has no literal
biclique with side sizes at least $\alpha n,\beta n$ with probability
$1-\exp(-\Omega(n^2))$. A nonempty literal biclique must have uniformly
positive literals on both sides or uniformly negative literals on both
sides. Each prescribed pair of sides of these sizes has probability at
most $2^{-\alpha\beta n^2}$ in either orientation. A union bound over
at most $3^n$ disjoint ordered pairs of variable sets and two orientations
proves the assertion.

Nevertheless, write $d_+(u)$ for the number of positive clauses incident
to $u$. In the absence of units,
\[
 F(u=1,y)-F(u=0,y)=d_+(u)-\sum_{v\ne u}y_v.
\]
Thus the combined comparison constraints, restricted to a common set of
free variables, are bounds on the same count. The one-sided constraints
from Section~\ref{sec:one} may have different literal-neighborhood
supports; they should not be described as identical constraints.

This family is a structural success case. If $m_{--}$ counts negative
clauses and $k=\sum_u x_u$, then
\begin{equation}
 F(x)=m_{--}+\sum_u d_+(u)x_u-\binom{k}{2}.
 \label{eq:sort}
\end{equation}
For each $k$, select the $k$ largest coefficients $d_+(u)$. Sorting once
and testing all counts solves the instance in polynomial time. Units only
modify the constant and the linear coefficients. The useful conclusion is
that redundant comparisons can reveal a cheaper optimization, rather than
that redundancy defeats the intended algorithm.

\section{Grouping variables by their effect across a cut}\label{sec:groups}

Use the Boolean quadratic representation
\begin{equation}
 F(x)=c+\sum_u h_u x_u+\sum_{u<v}J_{uv}x_ux_v,
 \qquad x_u\in\{0,1\}. \label{eq:quadratic}
\end{equation}
For one binary clause, its contribution to $J_{uv}$ is $-1$ for equal
literal polarities and $+1$ for opposite polarities. A missing clause
contributes zero. This is the Boolean coefficient convention; in the
spin convention $\sum h_u a_u-\sum J_{uv}^{\rm spin}a_ua_v$ of the
near-twin note, the coefficient on a simple edge is
$J_{uv}=-J_{uv}^{\rm spin}$. Unit clauses contribute only fields and a
constant. Fixed variable complementations may be performed first.

Let $A,B$ be disjoint and nonempty. Partition $A$ into nonempty groups
$A_1,\ldots,A_q$ such that
\begin{equation}
 J_{uv}=J_{u'v}\quad
 (u,u'\in A_i,\ v\in B). \label{eq:profiles}
\end{equation}
Write this common coefficient as $j_{iv}$. No restrictions are imposed
on interactions within $A$, within $B$, with the outside variables,
or on individual fields. In particular, the groups need not consist of
twins in the full graph.

Define
\begin{equation}
 P(A;B)=\prod_{i=1}^q(|A_i|+1),\qquad
 \Gamma(A;B)=|A|-\log_2 P(A;B)
 =\sum_{i=1}^q\bigl(|A_i|-\log_2(|A_i|+1)\bigr).
 \label{eq:gamma}
\end{equation}
For fixed literal orientations, the maximal identical-profile groups are
found by comparing the rows of $J_{A,B}$ in polynomial time. Merging two
groups of sizes $a,b$ that have the same profile increases $\Gamma$ by
$\log_2((a+1)(b+1)/(a+b+1))$. Thus using the maximal groups is optimal
among refinements of that partition for this bound.

\begin{theorem}[Group-count solver]\label{thm:groups}
Given $A,B$ and a partition satisfying~\eqref{eq:profiles}, the residual
Max 2-SAT instance on $n$ variables can be solved exactly in polynomial
space and time
\begin{align}
 \Ostar\!\left(2^{n-|A|-|B|}
       \left(2^{|A|}+P(A;B)2^{|B|}\right)\right)
 &=\Ostar\!\left(2^{n-|B|}+2^{n-\Gamma(A;B)}\right)\notag\\
 &=\Ostar\!\left(2^{n-\min\{|B|,\Gamma(A;B)\}}\right).
 \label{eq:group-time}
\end{align}
The number of groups is unrestricted.
\end{theorem}
\begin{proof}
Assign $R=V\setminus(A\cup B)$. For each count vector
$k=(k_1,\ldots,k_q)$ with $0\le k_i\le |A_i|$, the cross contribution is
\[
 \sum_i\sum_{u\in A_i}\sum_{v\in B}J_{uv}x_ux_v
 =\sum_{v\in B}\left(\sum_i k_i j_{iv}\right)x_v.
\]
The fixed assignment to $R$ contributes a constant and unary fields on
the two sides. After fixing $k$, the optimum therefore separates into
the optimum on $A$ with its group counts fixed and an unrestricted
optimum on $B$ with the displayed additional fields.

Generate count vectors one at a time. For the current vector, enumerate
only assignments to $A$ with these counts, retaining the best value and
witness, and enumerate the $2^{|B|}$ assignments to $B$. Across all count
vectors, the number of assignments to $A$ is exactly
\[
 \sum_{k_1=0}^{|A_1|}\cdots\sum_{k_q=0}^{|A_q|}
       \prod_i\binom{|A_i|}{k_i}=2^{|A|}.
\]
The second enumeration is repeated $P(A;B)$ times. This proves the first
bound, including all assignments to $R$. Formula~\eqref{eq:gamma}
gives the other expressions. Count vectors and group combinations can
be generated in polynomial space. Store only the current vector, best
assignments on the two sides for that vector, and a global best witness;
no table with $P(A;B)$ entries is required.
\end{proof}

\begin{corollary}[Linear repeated information]\label{cor:pairs}
Suppose $|B|\ge\beta n$ and $A$ contains $r\ge\eta n$ disjoint pairs
whose members have the same interaction profile on $B$. Their interactions
inside $A$ may be arbitrary. Then the instance is solvable in polynomial
space and time
\[
 \Ostar\!\left(2^{(1-\min\{\beta,(2-\log_2 3)\eta\})n}\right).
\]
\end{corollary}
\begin{proof}
Take the pairs as groups; any extra variables in $A$ can be singleton
groups. Each pair contributes $2-\log_2 3\simeq0.415$ to $\Gamma$;
singletons contribute zero. Apply Theorem~\ref{thm:groups}.
\end{proof}

More generally, if $|A|=\alpha n$, $|B|=\beta n$, and the number of
distinct profiles is $q=o(n)$, concavity gives
\[
 \log_2 P\le q\log_2(1+|A|/q)=o(n).
\]
The running time is then
$\Ostar(2^{(1-\min\{\alpha,\beta\})n+o(n)})$ in polynomial space.
This allows a growing number of groups, rather than only a bounded number
of interaction types.

\paragraph{Bicliques as one group.}
After orienting the literals, a complete literal biclique has
$J_{uv}=-1$ throughout the cross block. The same is true if every cross
clause is chosen arbitrarily between $(u\lor v)$ and
$(\neg u\lor\neg v)$. In both cases $q=1$,
$P=|A|+1$, and~\eqref{eq:group-time} recovers~\eqref{eq:biclique}.
Thus equal interaction profiles are strictly more permissive than equal
literal neighborhoods. Checking equality of profiles is a sufficient
structural test; logical redundancy among constraints has not been proved
to imply such equality.

\section{Branch on exceptions to expose repeated profiles}\label{sec:exceptions}

Let $X\subseteq B$. Keep $A$ fixed and regroup its variables by their
profiles on $B\setminus X$. Let $\Gamma_X$ be the resulting value
of~\eqref{eq:gamma}. Branch on every variable of $X$, then use the
group-count solver, also assigning $R$ as in its proof.

\begin{corollary}[Exception columns]\label{cor:exceptions}
This procedure runs in polynomial space and time
\begin{equation}
 \Ostar\!\left(2^{n-|B|+|X|}+2^{n-\Gamma_X}\right)
 =\Ostar\!\left(2^{n-S(X)}\right),\qquad
 S(X)=\min\{|B|-|X|,\Gamma_X\}.
 \label{eq:exceptions}
\end{equation}
\end{corollary}
\begin{proof}
There are $2^{|R|+|X|}$ assignments outside the two remaining sides.
The inner cost is
$2^{|A|}+P_X2^{|B|-|X|}$, where
$P_X=2^{|A|-\Gamma_X}$. Multiplying gives the two terms displayed.
Assigning $X$ changes fields, not the surviving cross profiles, so the
same partition works for every assignment to $X$.
\end{proof}

Deleting columns can only merge profile classes, so $\Gamma_X$ is
nondecreasing as $X$ grows. This gives a concrete branching score, with
both values of every exception variable already charged:
\begin{equation}
 \text{evaluate } S(X\cup\{v\})
 \quad\text{or the sharper sum in~\eqref{eq:exceptions}}
 \quad(v\in B\setminus X). \label{eq:selection}
\end{equation}
All one-column trials and their partitions are computable in polynomial
time. A selected variable is useful here because deleting its column
merges enough groups while leaving a sufficiently large opposite side.
This is a rule for making redundancy \emph{easier to exploit}, rather
than making literal imbalance grow.

For example, suppose $A$ contains linearly many pairs that disagree
arbitrarily within $A$, but each pair agrees on $B$ except possibly at
one common vertex $v$. Branching on $v$ exposes all these pairs at once.
If $B\setminus\{v\}$ is linear, Corollary~\ref{cor:pairs} gives a fixed
base below two. The earlier full-neighborhood twin rule need not apply
because of the arbitrary discrepancies inside $A$.

The score is a proved bound for a supplied cut and selected exception
set, not a guarantee that greedy deletion will discover the best set.
Several columns may need to be deleted before any classes merge.
Polynomial candidate lists or bounded lookahead are legitimate; an
exponential search for a promising partition must be included in the
running time. The symmetric construction, grouping $B$ by its effects
on $A$, is also available.

\section{A selective strategy and the desired theorem}\label{sec:target}

The following are components of a strategy, not a completed global
algorithm analysis:
\begin{enumerate}
\item Use domination and propagation, and consider both combined
comparison constraints and one-literal-neighborhood consequences.
Create or retain constraints when their incremental contribution to a
sound completion certificate justifies it. Constraints may also be
retained to exploit future shortening; no universal threshold on their
individual strength is imposed.
\item Choose exposure variables to improve the joint certificate,
shorten selected constraints, or improve the structural score
\eqref{eq:selection}. Assess all children, not only favorable outcomes.
\item When an identified cut has a large value of
$\min\{|B|,\Gamma(A;B)\}$, or of $S(X)$, replace further ordinary
branching by the exact group-count solver. Other certified structural
solvers, including the near-twin rule, remain available.
\end{enumerate}

For this terminal action, inherited pruning constraints can be discarded
and the residual \emph{soft} objective solved exactly. Every returned
assignment extends the currently fixed variables and is valid for the
original Max 2-SAT instance; the enlarged optimization domain contains
all retained completions. Its optimum supplies both a valid candidate
and an upper bound for the originally restricted subtree. Actual hard
constraints imposed as part of a different input problem could not be
discarded in this way.

Here is the quantitative target, specialized from the frontier formulation
of the learning-dichotomy note. It separates local structural savings from
the still-required bound on the number and cost of encountered states.

\begin{desiredtheorem}
For some density bound $\rho(n)\to\infty$ and constants
$\alpha,\gamma>0$ independent of $n$ and $\rho(n)$, there is a deterministic
polynomial-space branching strategy with polynomial work per initial
search node. Stop this initial search after at most
$\lfloor(1-\alpha)n\rfloor$ binary decisions on each path, or earlier at
a certified terminal solver. Its boundary consists of live nodes
$\mathcal F$ with continuation bounds $\Ostar(U_v)$ and structural
terminals $\mathcal S$ with solution-cost bounds $\Ostar(K_v)$, satisfying
\begin{equation}
 \sum_{v\in\mathcal F}U_v+\sum_{v\in\mathcal S}K_v
 \le 2^{(1-\gamma)n} \label{eq:frontier}
\end{equation}
up to polynomial factors. The initial target may restrict to one binary
clause per variable pair. The ultimate target covers arbitrary canonical
Max 2-SAT at density $O(n)$, including the two-clause encodings discussed
in the earlier notes.
\end{desiredtheorem}

At a group-count terminal with $n_v$ free variables, one admissible bound is
\[
 K_v=2^{n_v-|B_v|+|X_v|}+2^{n_v-\Gamma_{X_v}}.
\]
There are at most $O(2^{(1-\alpha)n})$ initial binary prefixes, and the
boundary continuations cost~\eqref{eq:frontier}. The desired theorem would
therefore give
$\Ostar(2^{(1-\min\{\alpha,\gamma\})n})$ time and polynomial space.
The terminal bound measures solution work, not the size of the feasible
domain. This permits cases such as~\eqref{eq:sort}, whose large domain
can be optimized without enumerating it.

The central remaining task is an inverse statement for states actually
reached by a specified strategy: when it cannot obtain enough additional
pruning, determine whether the objective has repeated profiles,
manageable discrepancies, count-based structure, or another certified
simplification. Absence of literal bicliques alone has not been shown
to imply this statement. Nor does failure of one chosen constraint or
one structural test refute it. A proof must quantify progress across
the entire frontier, rather than demand an immediate dichotomy at every
arbitrary collection of constraints.

\clearpage
\section{Relation to the earlier notes and verification}\label{sec:relation}

\paragraph{How this extends the structural rules.}
The interaction-class rule in
\path{max2sat_branching_frontier.tex} requires uniform interactions
inside and between its classes. The present theorem requires repetition
only across one chosen cut; interactions within both sides are arbitrary.
The near-twin rule in \path{max2sat_near_twin_pairs.tex} removes a pair
state after controlling discrepancies toward other simultaneously retained
pairs. Here the two sides are instead optimized separately after fixing
group counts. Neither rule must replace the other. The inherited
completion certificates and desired frontier budget are developed in
\path{max2sat_learning_dichotomy.tex}.

\paragraph{Reproducible checks.}
The companion program \path{explore_group_counts.py} compares the
streamed group-count algorithm with direct clause-by-clause exhaustive
optimization. It exercises arbitrary clauses within the sides and on
outside variables, missing cross edges, unequal unit biases, literal
complementations, and exceptional columns. It also checks the exact
flip identity and both necessary one-neighborhood bounds on every
assignment of small formulas. The output is
\path{group_count_checks.json}. With seed 87123, all 48 group-count
instances (including 24 with exception columns) and all 16 aligned
sorting instances matched exhaustive optimization. The verification
included 232,432 exact flip identities and 93,613 strict-improvement
cases satisfying both necessary one-neighborhood bounds. All streamed
assignment counts matched the two enumerations charged in the proof.
These checks test the implementation and
identities; the running-time and space bounds follow from the proofs,
and the desired theorem is not established by the experiments.

\end{document}
